\documentclass[11pt]{article}
\usepackage[margin=1in]{geometry}
\usepackage{graphicx}
\usepackage{booktabs}
\usepackage{hyperref}
\usepackage{amsmath}
\usepackage{float}
\usepackage{array}

\title{Searching for Tau Core signal candidates in SPARC residual structure}
\author{Jozsef Olcsak}
\date{2026}

\begin{document}
\maketitle

\begin{abstract}
Paper 1 found that externally reviewed structural disturbance is associated with larger low-acceleration residual scatter in SPARC. Paper 2 reversed the question and showed that fixed residual-shape features can recover those A/C labels better than chance, while remaining explicitly non-unique with respect to MOND-simple and empirical RAR-like baselines. This Paper 3 seed asks a narrower follow-up question: where do TPG/projection residuals, MOND-simple residuals, RAR-like residuals, and measured rotation curves point to candidate residual structure that could plausibly carry the missing observer- and environment-dependent Tau Core weight? The current packet identifies candidate galaxies, radial onset categories, and environment/observability stress channels. It does not validate Tau Core, does not numerically test RMOND, and does not claim gravity-model selection.
\end{abstract}

\section{Purpose and claim boundary}

The working hypothesis is that the TPG/projection baseline already carries part of the local Tau Core weighting, while the remaining TPG residual may carry missing observer- and environment-dependent weights. This is a candidate-search statement, not a proof. A defensible Paper 3 must therefore separate three layers:
\begin{enumerate}
\item an operational residual layer, based on frozen TPG/projection, MOND-simple, and RAR-like residual maps;
\item a systematics layer, based on distance, inclination, point count, H\,I kinematics, beam smearing, and non-circular motions;
\item a theory layer, where a Tau Core interpretation is allowed only if the residual pattern survives the systematics layer and is more specific than ordinary RAR/MOND behavior.
\end{enumerate}

The permitted claim is that this packet defines reproducible Tau Core signal candidates. The forbidden claim is that the candidates already prove Tau Core.

\section{Data inherited from Papers 1 and 2}

The packet inherits derived SPARC residual artifacts from the Paper 1 and Paper 2 public packages. Raw SPARC rotmod files and raw H\,I survey products are not redistributed. The retained derived tables include pointwise absolute residuals for TPG/projection, MOND-simple, and empirical RAR-like families; galaxy-level residual features; distance, observability, and environment summaries; and small radial-pilot control summaries.

The main residual quantity is
\[
r_{m,gi}=\left|\log V_{{\rm obs},gi}-\log V_{m,gi}\right|,
\]
with galaxy-level burden
\[
{\rm RMS}_{m,g}=\sqrt{{1\over N_g}\sum_i r_{m,gi}^2}.
\]
The TPG-specific excess used for candidate triage is
\[
\Delta_{{\rm TPG-low},g}=
{\rm RMS}_{{\rm TPG},g}-
\min\left({\rm RMS}_{{\rm MOND},g},{\rm RMS}_{{\rm RAR},g}\right).
\]
Positive values mark galaxies where the TPG/projection residual is larger than the best available low-acceleration comparator in this packet. Negative values mark TPG success controls.

\section{Theory motivation}

The relevant external baseline is the SPARC radial acceleration relation \cite{McGaugh2016RAR,Lelli2016SPARC} and the broader MOND/RAR literature \cite{Milgrom1983MOND,Li2018RARFits}. These works are not optional background; they define the strongest ordinary low-acceleration competitor. If Tau Core only reproduces the same residual ordering as RAR/MOND, it is not yet distinguished.

The systematics literature is equally central. Non-circular motions in H\,I velocity fields \cite{Trachternach2008THINGS,Oman2019NonCircular} and beam-smearing/rotation-curve quality effects \cite{deBlok1997Beam} can generate residual structure without new physics. The JWST/NIRCam Zone-of-Avoidance result \cite{NiloCastellon2025ZoA} is used only as an observer/line-of-sight motivation: foreground obscuration and hidden structure can materially affect what an observer can map. It is not evidence for Tau Core by itself.

\section{Candidate construction}

For each galaxy we record TPG/projection RMS, MOND-simple RMS, RAR-like RMS, TPG-specific excess, environment proxy, distance, reconstruction-risk proxy, inclination, and the first radial bin where the TPG/projection absolute residual exceeds 0.15 dex. Candidate classes are assigned by frozen screening rules:
\begin{itemize}
\item \textit{clean tau candidate}: high TPG residual, high TPG-specific excess, and non-high reconstruction risk;
\item \textit{environment tau candidate}: high residual burden with high environment cue;
\item \textit{TPG divergence follow-up}: positive TPG-specific excess without enough cleanliness for the first two labels;
\item \textit{TPG success control}: low residual burden where TPG is no worse than MOND/RAR-like baselines.
\end{itemize}

These classes are triage labels. They are not physical classifications.

\begin{table}[H]
\centering
\caption{Candidate shortlist from the regenerated seed packet.}
\scriptsize
\setlength{\tabcolsep}{3pt}
\begin{tabular}{llrrrp{0.22\linewidth}p{0.18\linewidth}}
\toprule
Galaxy & Class & TPG RMS & TPG excess & Env. & Onset & Candidate class\\
\midrule
UGC00731 & C & 0.340 & 0.093 & 0.00 & inner\_R$<$0.33Rmax & clean tau \\
CamB & A & 0.964 & -0.080 & 1.72 & inner\_R$<$0.33Rmax & environment \\
UGC05764 & A & 0.329 & 0.076 & 0.00 & inner\_R$<$0.33Rmax & TPG divergence \\
NGC2403 & A & 0.085 & 0.029 & 1.86 & outer\_R$>$=0.67Rmax & TPG divergence \\
NGC3741 & C & 0.204 & 0.103 & -0.12 & inner\_R$<$0.33Rmax & TPG divergence \\
UGC08490 & C & 0.186 & 0.055 & 0.00 & inner\_R$<$0.33Rmax & TPG divergence \\
NGC3109 & C & 0.176 & 0.039 & 0.20 & middle\_0.33-0.67Rmax & TPG divergence \\
UGC06446 & C & 0.194 & 0.038 & 0.00 & inner\_R$<$0.33Rmax & TPG divergence \\
\bottomrule
\end{tabular}
\end{table}

\section{Residual and environment readouts}

\begin{figure}[H]
\centering
\includegraphics[width=0.82\linewidth]{figures/paper3_environment_residual_scatter.pdf}
\caption{TPG/projection residual burden versus the inherited environment proxy. Named points are the largest seed-candidate scores.}
\end{figure}

\begin{figure}[H]
\centering
\includegraphics[width=0.82\linewidth]{figures/paper3_tpg_specific_excess.pdf}
\caption{Galaxies with the largest TPG/projection residual excess relative to the better of MOND-simple and RAR-like RMS.}
\end{figure}

\begin{figure}[H]
\centering
\includegraphics[width=0.78\linewidth]{figures/paper3_residual_onset_bins.pdf}
\caption{Radial bin where the TPG/projection absolute residual first exceeds 0.15 dex. This is a morphology-of-failure diagnostic, not a detection statistic.}
\end{figure}

\begin{table}[H]
\centering
\caption{Screening correlations for TPG/projection residual burden.}
\footnotesize
\begin{tabular}{lrl}
\toprule
Covariate & Pearson $r$ & Interpretation\\
\midrule
DistanceMpc & -0.291 & screening\_correlation\_not\_causal\_attribution \\
EnvMaxTheta & 0.110 & screening\_correlation\_not\_causal\_attribution \\
W\_tau\_eff\_abs\_v01 & 0.983 & screening\_correlation\_not\_causal\_attribution \\
ReconstructionRiskChannel\_v01 & 0.317 & screening\_correlation\_not\_causal\_attribution \\
MeanErrVobsKms & -0.214 & screening\_correlation\_not\_causal\_attribution \\
InclinationErrorDeg & 0.214 & screening\_correlation\_not\_causal\_attribution \\
\bottomrule
\end{tabular}
\end{table}

\begin{figure}[H]
\centering
\includegraphics[width=0.82\linewidth]{figures/paper3_observability_stress.pdf}
\caption{Screening correlations with the TPG/projection residual burden. These are not causal estimates.}
\end{figure}

\section{Required $S_\tau$ diagnostic}

The extended Tau Core form can be written operationally as
\[
F_\tau(a_N,R)=1+S_\tau(R)\,\alpha\ln\left(1+{a_0\over a_N}\right).
\]
For each SPARC point with usable baryonic baseline, the required local value is
\[
S_{\tau,\mathrm{req}}(R)=
{V_{\rm obs}(R)/V_N(R)-1\over \alpha\ln\left(1+a_0/a_N(R)\right)}.
\]
This is an inverse diagnostic. It uses the measured endpoint and therefore cannot be used as a predictive Tau Core validation. Its role is to ask what kind of future predictive rule would be needed: a galaxy-level constant, a radial function, an acceleration function, or an environment-coupled function.

\begin{figure}[H]
\centering
\includegraphics[width=0.78\linewidth]{figures/paper3_required_s_tau_scatter.pdf}
\caption{Median required $S_\tau$ versus within-galaxy spread. A low spread near $S_\tau=1$ would support a constant correction; large spread indicates that a function is needed.}
\end{figure}

\begin{table}[H]
\centering
\caption{In-sample required-$S_\tau$ function-family diagnostics. Lower RMS means that family can absorb the measured TPG residual more compactly.}
\footnotesize
\begin{tabular}{lrrr}
\toprule
Family & Median RMS & Mean RMS & Galaxies\\
\midrule
radius + acceleration & 0.143 & 0.152 & 45 \\
quadratic radius & 0.145 & 0.153 & 45 \\
linear radius & 0.146 & 0.168 & 45 \\
quadratic acceleration & 0.165 & 0.160 & 45 \\
linear acceleration & 0.165 & 0.167 & 45 \\
galaxy constant & 0.187 & 0.183 & 45 \\
fixed S=1 & 0.188 & 0.249 & 45 \\
\bottomrule
\end{tabular}
\end{table}

\begin{figure}[H]
\centering
\includegraphics[width=0.82\linewidth]{figures/paper3_s_tau_family_comparison.pdf}
\caption{Diagnostic comparison of constant and simple functional $S_\tau$ families. This is not model selection because the measured velocities are used to infer $S_\tau$.}
\end{figure}

\section{RMOND and comparator status}

The requested RMOND comparison is not yet a numeric endpoint in this seed packet. The local TPG--RMOND bridge is useful, but it is not the same object as a frozen pointwise velocity law. Its strongest result is structural compatibility: the tau-projection can preserve the leading SZ20 scalar and vector Lagrangian exponents under a channel-normalised kernel. That is theory motivation, not a direct prediction for $V(R)$.

The local RMOND--MTW hybrid analysis also warns against double counting. At $f_{\rm vec}=0$, the hybrid reproduces pure MTW/TPG exactly, so it is not an independent comparator. At nonzero $f_{\rm vec}$, the vector contribution requires a frozen halo/vector prescription and screening rule before any residual table is meaningful. Treating either option as a completed RMOND residual test would weaken the paper.

The next technical gate is therefore sharper than ``find a CSV'': freeze a unique $V_{\rm RMOND}(R)$ prescription, regenerate a derived pointwise residual table on the same SPARC radii, and only then add RMOND beside TPG/projection, MOND-simple, and RAR-like baselines.

\section{Interpretation}

The most useful interpretation is modest. TPG success controls show where the local projection baseline is already adequate. TPG divergence follow-ups show where the residual has structure that might encode missing observer/environment weights, but these are exactly the cases where ordinary systematics can also enter. A candidate becomes interesting only when three facts hold together: the TPG residual diverges in a structured radial way, the divergence is not shared by all low-acceleration baselines, and the object has a plausible Tau Core weighting channel that is not reducible to distance, inclination, beam smearing, or non-circular motion.

\section{Conclusion}

This seed opens Paper 3 as a reproducible candidate-search project. It carries forward the discipline learned from Papers 1 and 2: freeze the endpoint, compare against MOND/RAR baselines, name the systematics, and avoid promoting a diagnostic pattern into a physical proof. The current evidence is enough to define targets for Tau Core follow-up, especially TPG-divergence galaxies and TPG-success controls, but it is not enough to claim a detected Tau Core field.

The next paper-grade step is a frozen RMOND residual table plus a held-out environment/line-of-sight validation rule. Only after that can the paper ask whether the residual candidates are truly Tau Core-specific rather than ordinary low-acceleration or H\,I-systematics behavior.

\bibliographystyle{plain}
\bibliography{references}

\end{document}
